% Zmiennosc wartosci koncowych ABC vs Bee (boxplot + surowe punkty).
% Dane: grafiki/final_values.dat (generowane przez: python experiments/run_experiments.py)
\begin{tikzpicture}
\begin{axis}[
    wykres,
    width=0.55\textwidth,
    height=7.2cm,
    boxplot/draw direction=y,
    xtick={1,2},
    xticklabels={ABC, Bee},
    xmin=0.4, xmax=2.6,
    ylabel={Końcowa wartość rozwiązania},
    xmajorgrids=false,
]
% mark=none: surowe punkty rysujemy osobno, bez podwojnego znacznika outliera
\addplot[boxplot={draw position=1, box extend=0.45}, draw=abcblue, fill=abcblue!20, solid, mark=none]
    table [y=abc] {grafiki/final_values.dat};
\addplot[boxplot={draw position=2, box extend=0.45}, draw=beeorange, fill=beeorange!20, solid, mark=none]
    table [y=bee] {grafiki/final_values.dat};
% surowe wyniki pojedynczych uruchomien
\addplot[only marks, mark=*, mark size=1.5pt, mark options={fill=black!70, draw=black!70}]
    table [x expr=1, y=abc] {grafiki/final_values.dat};
\addplot[only marks, mark=*, mark size=1.5pt, mark options={fill=black!70, draw=black!70}]
    table [x expr=2, y=bee] {grafiki/final_values.dat};
\end{axis}
\end{tikzpicture}
